\section{Discussion: when extra loops at inference help}
\label{sec:discussion}

The same architecture is parameter-efficient on text and, on the chain task, carries out a serial computation that more
loops can extend. Our interpretation of the difference, which we have not tested, is that extra loops at inference help when the model's
errors come from too few serial steps and the loop can run longer without destroying its state, and that using them in
deployment also needs a signal for when to stop.

The first condition is the hardest to measure on its own, because here it is inferred from where extra loops help. On
the chain task, the dense models know every table, and within the training range their errors come from too few serial
steps: one more loop is one more step. On web text at these sizes we found no sign of such errors. The per-token gain is
spread across tokens, and extra loops never helped, at any learning rate. The larger gains on math and code text hint at
more errors of this kind there, but at this size the results on those sets vary mainly with the initialisation and the
data order.

Whether a loop can run longer depends on the kind of iteration that training gives it. On text, a middle block with a
fixed loop count becomes a fixed-length program, a whole stack trained with truncated backpropagation a periodic cycle,
and a middle block with a random loop count a fixed point. The fixed point can be run for any number of loops but gains nothing on text from the extra ones, the cycle works only
at multiples of its period, and the fixed-length program works only at its training count. On the chain task, models trained with as many loops as steps lose their values
when run longer, and a few random extra loops in training fix that.

To spend extra loops only where a problem needs them, a deployed model must decide how many to run. A fixed number of
extra loops works for the 8-loop model but not for the loops = steps model, which loses its values when run longer. On
the chain task, the model's own confidence served as such a signal for most training regimes \lbl{(exploratory)}, including the one
that extends furthest once it is trained with a few loops past the last step. We
did not test stopping rules on text, where extra loops did not help.

The chain task meets these conditions once the model is trained for them. Web text at these sizes showed no sign of
the first, and the one text regime that can run longer, the random loop count, gained nothing from it. Two follow-ups
that we did not run would test this interpretation. One is a serial task whose input length does not grow with the
number of steps, or this task with prompts padded to a fixed length. It would also show whether a written chain of
thought extends when the prompt does not get longer, which is the control that Section~\ref{sec:extension} lacks and
that no frozen plan included. The other is training on text that needs more serial steps, such as mathematics and code,
compared at equal compute.

On text, where extra loops do not help, a loop still has value at equal parameters, and how much depends on size, data
and the dense model it is compared with (Section~\ref{sec:text}). Some effective-parameter laws let the factor applied to
the looped parameters depend on the loop count but not on model size or training tokens. Such laws fit the multiplier of
the larger models but not its rise with data, and they do not fit the smaller models at all; the single pre-registered
$\varphi$ is an average over these two regimes. Studies at one scale, or against dense references of one depth, will
underestimate or overestimate the loop, depending on their scale and reference.

\paragraph{Practical guidance.} Within the range we measured:
\begin{itemize}
  \item On text, use a loop when parameters are the binding constraint, rather than data, compute or latency. Its
  advantage at equal parameters grows with model size and tokens per parameter, but a dense model was better at equal
  training compute in every measured comparison but one (Section~\ref{sec:equalcompute}), and at deployment each
  executed layer costs decoding time and KV cache.
  \item Try a deeper dense model of the same size first. At 40M it gave about half of the advantage for half the
  training compute, and at the shortest budget (10 tokens per parameter), with its learning rate tuned, all of it (one
  seed). At 20M, at width 320, it was worse than the 8-layer model.
  \item If you loop, loop a middle block with full backpropagation. At similar executed depth, the middle block won
  every comparison with the whole stack (Table~\ref{tab:verdicts}). Start with four loops on text; at 40M--80M, eight did
  about as well at about 1.75 times the training compute. Judge at the target scale, since small-scale ablations
  underestimate the loop.
  \item Tune the looped model's learning rate separately, especially at small scale and on tasks that need serial steps.
  At a rate that suited the dense model, most looped runs did not learn the chain task within the given training steps.
  \item Do not expect extra loops at inference to help on text after standard training. Sweep the loop count at
  inference to see whether a model can use more loops, and probe it to see whether it computes step by step even if it
  cannot.
  \item If a model must decide how long to run, its confidence was the most useful stopping signal here
  \lbl{(exploratory)}. Training a few loops past the last step made it work for the loops = steps model in all five
  seeds. Within the training range of our task, a written chain of thought, supervised on every intermediate value, was
  the cheaper choice; its failure beyond the range probably came from the longer prompts and says little about inputs of
  fixed length.
\end{itemize}
